% Part: computability
% Chapter: computability-theory
% Section: rice-theorem

\documentclass[../../../include/open-logic-section]{subfiles}

\begin{document}

\olfileid[fr]{cmp}{thy}{rce}
\olsection{Le théorème de Rice}

Les particularités de $\fn{Tot}$ ne jouent aucun
rôle dans la preuve de la \olref[tot]{prop:total}.
Nous avons conçu $h(x,y)$ pour qu'elle coïncide
avec la fonction constante $j(y)=0$ exactement
lorsque $x\in K$; nous aurions pu la faire
coïncider, dans ce cas, avec n'importe quelle
autre fonction calculable partielle. Cette
observation mène à un théorème plus général :
grossièrement, aucune propriété non triviale
des fonctions calculables partielles n'est décidable.

Rappelons que $\cfind{0}$, $\cfind{1}$,
$\cfind{2}$,~\dots est notre énumération
canonique des fonctions calculables partielles.

\begin{thm}[Théorème de Rice]
Soit $C$ un ensemble quelconque de fonctions
calculables partielles, et posons
$A=\Setabs{n}{\cfind{n}\in C}$.
Si $A$ est calculable, alors soit $C$ est vide,
soit $C$ est l'ensemble de toutes les fonctions
calculables partielles.
\end{thm}

Un \emph{ensemble d'indices} est un ensemble~$A$
tel que, si les indices $n$ et~$m$ calculent la
même fonction, ils appartiennent tous deux à~$A$
ou n'y appartiennent ni l'un ni l'autre.
L'ensemble~$A$ du théorème possède manifestement
cette propriété. Réciproquement, si $A$ est un
ensemble d'indices et $C$ l'ensemble des fonctions
calculées par ces indices, alors
$A=\Setabs{n}{\cfind{n}\in C}$.

\begin{explain}
En ces termes, le théorème de Rice affirme
qu'aucun ensemble d'indices non trivial n'est
décidable. Pour en comprendre la portée, il
faut distinguer les \emph{programmes}, objets
syntaxiques d'un langage de programmation, des
fonctions qu'ils calculent. Certaines questions
sur les programmes et leurs indices sont
effectivement décidables : ce programme contient-il
plus de 150 symboles ? Plus de 22 lignes ?
Une instruction « while » ? La chaîne
« hello world » apparaît-elle comme argument
d'une instruction « print » ? Le théorème de
Rice dit qu'aucune question non triviale sur la
\emph{fonction} calculée par un programme
n'est décidable. Par exemple : le programme
s'arrête-t-il sur l'entrée~$0$ ? S'arrête-t-il
pour au moins une entrée ? Produit-il parfois
un nombre pair ?
\end{explain}

\begin{proof}[Démonstration du théorème de Rice]
Supposons que $C$ ne soit ni vide ni égal à
l'ensemble de toutes les fonctions calculables
partielles, et que $A$ soit l'ensemble des indices
des fonctions de~$C$. Nous montrerons que, si
$A$ était calculable, on pourrait résoudre le
problème de l'arrêt. Donc $A$ n'est pas calculable.

On peut supposer, sans perte de généralité,
que la fonction~$f$ nulle part définie n'appartient
pas à~$C$; dans le cas contraire, on échange~$C$
et son complémentaire parmi les fonctions
calculables partielles. Soit $g$ une fonction
quelconque de~$C$. Si l'on pouvait décider~$A$,
on pourrait distinguer les indices calculant~$f$
de ceux calculant~$g$, et exploiter cette
distinction pour résoudre le problème de l'arrêt.

Voici comment. À l'aide de la fonction partielle
calculable universelle, définissons
\[
h(x,y) \simeq
\begin{cases}
\text{indéfinie} & \text{si $\cfind{x}(x) \fundefined$} \\
g(y) & \text{sinon.}
\end{cases}
\]
Pour calculer $h$, on lance d'abord le calcul
de $\cfind{x}(x)$. S'il s'arrête, on calcule
ensuite $g(y)$ et, si ce second calcul s'arrête
à son tour, on renvoie sa valeur. Plus
formellement, on peut écrire
\[
h(x,y) \simeq \Proj{2}{0}(g(y),\fn{Un}(x,x)).
\]
Ici $\Proj{2}{0}(z_0,z_1)=z_0$ est la projection
calculable à deux arguments qui renvoie le
premier. La fonction~$h$ est donc une composée
de fonctions calculables partielles : le membre
de droite est défini et vaut~$g(y)$ exactement
lorsque $\fn{Un}(x,x)$ et $g(y)$ sont tous deux
définis.

Pour $x$ fixé, si $\cfind{x}(x)$ n'est pas
définie, $h(x,y)$ n'est définie pour aucun~$y$;
si $\cfind{x}(x)$ est définie, alors
$h(x,y)\simeq g(y)$. Ainsi, selon~$x$,
la fonction $y\mapsto h(x,y)$ coïncide soit
avec~$f$, soit avec~$g$. Décider si
$\cfind{x}(x)$ est définie revient à savoir
laquelle de ces deux possibilités se réalise,
ou encore si la fonction
$h_x(y)\simeq h(x,y)$ appartient à~$C$.
Un test calculable d'appartenance à~$A$
permettrait de le savoir.

Plus précisément, $h$ étant calculable partielle,
elle possède un indice~$e$ tel que
$h\simeq\cfind{e}[2]$. Le théorème
$s$-$m$-$n$ fournit une fonction récursive
primitive~$s$ telle que, pour chaque~$x$,
$\cfind{s(e,x)}(y)\simeq h_x(y)$.
Si $\cfind{x}(x)\fdefined$, alors
$\cfind{s(e,x)}$ coïncide avec~$g$ et
$s(e,x)\in A$. En revanche, si
$\cfind{x}(x)\uparrow$, elle coïncide
avec~$f$ et $s(e,x)\notin A$. Pour tout~$x$,
on a donc $x\in K$ si et seulement si
$s(e,x)\in A$. Si $A$ était calculable,
$K$ le serait aussi, contradiction.\footnote{Corrections éditoriales : $h$ est binaire, donc son indice s'écrit $h\simeq\cfind{e}[2]$, plutôt que $h=\cfind{e}$ dans la source. Nous employons $\simeq$ dans la spécialisation. La portée du théorème est la fonction partielle calculée, non toute propriété du comportement d'exécution.}
\end{proof}

Le théorème de Rice a de nombreuses applications.
Le corollaire suivant en donne quelques exemples.
\begin{cor}
Les ensembles suivants sont indécidables :
\begin{enumerate}
\item $\Setabs{x}{\text{$17$ appartient à l'image de $\cfind{x}$}}$;
\item $\Setabs{x}{\text{$\cfind{x}$ est constante}}$;
\item $\Setabs{x}{\text{$\cfind{x}$ est totale}}$;
\item $\Setabs{x}{\text{pour tous $y<y'$, $\cfind{x}(y) \fdefined$ et,
    si $\cfind{x}(y') \fdefined$, alors $\cfind{x}(y) < \cfind{x}(y')$}}$.
\end{enumerate}
\end{cor}

\begin{proof}
Ce sont tous des ensembles d'indices non triviaux. Dans le quatrième cas, prendre $y'=y+1$ montre que la condition impose la totalité ; elle caractérise donc les fonctions totales strictement croissantes, et non les seules fonctions croissantes sur leur domaine partiel.
\end{proof}

\end{document}
